\begin{textsample}{POS Dim 1 – mt – Score 48.00 – mt/mt\_ef\_71.txt}  \label{ex:f1_pos_077}
2.2.2.
Contribuição de Ciências para o Desenvolvimento Integral Os conceitos e procedimentos da área de Ciências da Natureza colaboram para a compreensão do mundo e suas transformações, porque contribuem para o entendimento e o questionamento dos diferentes modos de intervenção e para o reconhecimento do ser humano como parte integrante do Universo.
Nessa perspectiva, Sasseron ( 2011 ) traz \textbf{que}, devido à pluralidade semântica, encontra -se \textbf{hoje} em dia, na literatura nacional sobre ensino de Ciências, \textbf{autores} \textbf{que} utilizam a expressão “ Letramento Científico ” ( Mamede e Zimmermann, 2007, Santos e Mortimer, 2001 ).
Outros \textbf{pesquisadores} adotam o termo “ Alfabetização Científica ” ( Brandi e Gurgel, 2002, Auler e Delizoicov, 2001, Lorenzetti e Delizoicov, 2001, Chassot, 2000 ) e também aqueles \textbf{que} usam a expressão “ Enculturação Científica ” ( Carvalho e Tinoco, 2006, Mortimer e Machado, 1996 ) para designarem o objetivo desse ensino de Ciências \textbf{que} almeja a formação cidadã dos estudantes para o domínio e uso dos conhecimentos científicos e seus \textbf{desdobramentos} nas mais diferentes esferas de sua vida.
Portanto, percebe -se \textbf{que} alguns \textbf{pesquisadores} \textbf{que} usam \textbf{um} termo ou outro partilham das mesmas premissas no ensino de Ciências, ou seja, os motivos \textbf{que} guiam o planejamento desse ensino para a construção de benefícios práticos para as pessoas, a sociedade e o ambiente.
Assim, a alfabetização e o letramento científico tornaram -se mais \textbf{que} \textbf{uma} necessidade, \textbf{uma} exigência, pois permite ampliar a forma de \textbf{ver}, sentir, apreciar e fazer uso dos recursos naturais, sobretudo, \textbf{que} permita ao estudante identificar o alcance da aplicação científica entre jovens e adultos em seu cotidiano, de maneira consciente e sustentável, a fim de estabelecer relações de harmonia entre Ciência, Sociedade e Tecnologia.
Nessa perspectiva, \textbf{uma} educação inovadora em Ciências, contextualizada, \textbf{que} gere indagações e interesse pelas descobertas, \textbf{que} construa e ofereça respostas \textbf{conceituais} \textbf{que} contribuam para a melhoria da qualidade de vida, faz -se necessária.
Nessa linha de \textbf{raciocínio}, o letramento científico é \textbf{um} diferencial para o entendimento da ciência e sua utilização pela sociedade em \textbf{tarefas} cotidianas, como, por exemplo, a leitura dos dados em \textbf{uma} conta de energia, compreensão de \textbf{uma} bula de medicamentos, leitura da conta de água e análise do volume consumido.
Dessa maneira, torna -se latente a \textbf{percepção} da importância de se dar \textbf{um} tratamento prático aos conteúdos da \textbf{disciplina} de Ciências.
Esse é o desafio para a prática docente, proporcionar aos estudantes a compreensão da relevância científico-tecnológica no cotidiano e no processo de entendimento dos fenômenos naturais, por meio de atividades \textbf{que} privilegiam vivência e reflexões das ações diárias, de forma a contextualizar os conteúdos \textbf{inerentes} aos eixos temáticos previstos nos PCN para o Ensino Básico ( BRASIL, 1997 ) e às Unidades Temáticas definidas na Base Nacional Comum Curricular ( BRASIL, 2017 ).
Assim, entendendo \textbf{que} não é possível mudar a natureza humana pelo fato de estar em espaços culturais diversos, é necessário ajudar os estudantes a interpretarem suas realidades em diferentes contextos, primando pelo respeito às diversidades, sejam elas culturais, sociais, linguísticas ou de qualquer outra ordem.
Diante desse entendimento, o ensino de Ciências, no ensino fundamental, representa a base para a educação científica e entendimento de mundo, com intenção, não apenas de garantir \textbf{que} a aprendizagem tenha significado para o estudante, mas \textbf{que}, além disso, contribua para sua autonomia intelectual, \textbf{que} \textbf{implique} em ações voltadas para sustentabilidade \textbf{ancorada} na ética, na cidadania e na democracia ( BNCC, 2017 ).
A apresentação dos objetos de conhecimento voltados ao ensino-aprendizagem, durante o ano letivo, não devem ser ministrados somente por meio de aulas expositivas, pois nos parece \textbf{um} caminho mais \textbf{interessante} ( ou, ao menos alternativo ) a utilização de estratégias lúdicas \textbf{que} utilizam vídeos, softwares para simulação de eventos científicos, atividades extraclasses, projetos \textbf{que} incluam a internet das coisas, por exemplo, na busca por soluções para situações-problemas relacionadas com transporte público, demanda de energia, sustentabilidade, clima, reservas de água e consumo, produção de alimentos, saúde pública e, por fim, alternativas \textbf{que} privilegiem o protagonismo do estudante no ensino-aprendizagem de ciências.
A realização de atividades diferenciadas valoriza a participação ativa dos estudantes no sentido de promover a interação, por meio da linguagem, para a formação de novos conceitos prévios na estrutura cognitiva.
O desenvolvimento da \textbf{mente} e do corpo ocorre de maneira integrada e não dissociada do mundo \textbf{que} rodeia a todos, entendida nesse contexto, a vida recebe \textbf{influências} do meio em \textbf{que} as pessoas \textbf{vivem}, dos alimentos \textbf{que} ingerem e das mudanças sociais, econômicas, políticas, tecnológicas e ambientais \textbf{que} ocorrem o tempo todo, até porque o ser humano não é estanque, seu movimento está em constante modificação, interação, evolução, integração e adaptação.
O ensino de Ciências, nessa etapa de redescoberta, deve ser o diferencial \textbf{que} irá gerar o desejo de conhecer e, por \textbf{conseguinte}, \textbf{despertar} a sensação de deslumbramento pela descoberta, proporcionará \textbf{um} ensino-aprendizagem \textbf{que} aponte o \textbf{lado} filosófico da ciência, favorecendo o letramento científico numa perspectiva mais ampla, em \textbf{que} o estudante \textbf{aprimore} a capacidade de empregar o conhecimento científico para reconhecer questões, entender e explicar fenômenos científicos, realizar conclusões \textbf{baseadas} na análise do agrupamento de dados coletados na pesquisa.
Nesse sentido, é \textbf{inerente} ao conceito de letramento científico a concepção e diferenciação das características indutoras de \textbf{que} a ciência é \textbf{uma} forma de conhecimento e investigação.
Assim, o estudante deve ser capaz de perceber \textbf{que} a ciência e a tecnologia moldam o meio material, cultural e intelectual ; e, por fim, sensibilizar como cidadão crítico e preocupado em compreender, analisar tomar decisões sobre as transformações ocorridas.
A formação científica acontece quando a curiosidade, a observação e a experimentação são associadas à visão de mundo, ampliando e aprofundando os conhecimentos prévios.
Os \textbf{fundamentos} pedagógicos da Base nacional Comum Curricular, homologada em dezembro de 2017, demonstram preocupação com o desenvolvimento global ( dimensões intelectual, física, afetiva, social, ética, moral e simbólica ) e \textbf{compromisso} com a educação integral ( independente da duração da jornada escolar ), apontando para a construção de processos educativos sintonizados com as necessidades, possibilidades e interesse dos estudantes, sem desconsiderar, \textbf{contudo}, os desafios da sociedade \textbf{contemporânea} para formar pessoas autônomas e capazes de se servir dessas aprendizagens ao longo da vida.
Na BNCC, objetos de conhecimento estão a serviço das competências ( cognitivas, comunicativas, pessoais e sociais ) \textbf{que} visam à formação humana integral e à construção de \textbf{uma} sociedade \textbf{justa}, democrática e inclusiva.
As competências gerais são a base para as competências das áreas \textbf{que}, por sua vez, \textbf{ancoram} as competências dos componentes curriculares.
A Área de Ciências da Natureza é composta pela \textbf{disciplina} de Ciências \textbf{que} se divide em três Unidades Temáticas.
A partir das transformações da matéria/energia, da degradação ambiental e suas implicações para a manutenção da vida, aborda -se a Unidade temática Matéria e Energia, \textbf{uma} vez \textbf{que} esses assuntos contribuem para o entendimento da organização humana e suas relações de sustentabilidade com o meio.
Essa \textbf{abordagem} \textbf{permeia} o tempo presente e o passado em diferentes contextos, além de considerar as questões éticas, valores socioambientais e ações coletivas coerentes com o bem comum.
Assim, por meio da Unidade Temática Vida e Evolução promove -se a ampliação do conhecimento acerca da diversidade da vida nos ambientes naturais ou transformados pelo ser humano, a compreensão da integridade do corpo através do tema alimentação, por exemplo, e como ela favorece ou não a saúde e as doenças, abordando o assunto com \textbf{interface} nos diferentes hábitos alimentares culturais, estabelecendo relações entre os processos vitais da fisiologia humana.
Por fim, por \textbf{intermédio} da Unidade temática Terra e Universo é possível explorar aspectos relacionados à manutenção da vida no planeta, processos biogeoquímicos \textbf{que} favorecem a perpetuação das espécies, bem como a explicação de fenômenos relacionados aos astros e à \textbf{influência} das várias definições \textbf{que} ocorreram ao longo da história por diferentes povos e culturas.
O processo de globalização conferiu novas realidades à escola e à educação, na acepção de instituição responsável pelo ensino formal, entendendo -a, não apenas como disseminadora de informações tidas como privilegiadas.
Nesse contexto, não pode ser apenas transmissora de conteúdos, \textbf{teorias} ou conceitos científicos, deve potencializar alternativas \textbf{que} privilegiam o entendimento do eu, do outro e das relações destes com o universo e suas manifestações, para \textbf{que} o estudante tenha condições de, a partir desses pressupostos, entender o mundo natural, propondo ações \textbf{que} conduzam a \textbf{uma} qualidade de vida cada vez melhor.
Para Chassot ( 2003 ), “ propiciar o entendimento ou a leitura dessa linguagem é fazer alfabetização científica ”.
A Internet das Coisas Para o Ensino de Ciências da Natureza e Tecnologias sugere -se a IoT como \textbf{uma} ferramenta facilitadora de ensino aprendizagem, assim como mapas \textbf{conceituais}, as atividades lúdicas, os simuladores de \textbf{ensino-aprendizagem} e os kits de ensino de Ciências da Natureza.
No \textbf{que} se refere à evolução da Internet, \textbf{conhecida} como IoT ( Internet of Things ), também chamada de Internet dos objetos, é \textbf{uma} nova referência na história, levando em consideração o impacto \textbf{que} já teve na educação, proporcionando a pesquisa e a divulgação de seus resultados, contribuindo para a simulação de fenômenos e análise, \textbf{enquanto} divulgação do \textbf{que} é ciências, de acordo com sua epistemologia, critérios de validação, na comunicação, nos negócios, no governo e na \textbf{humanidade}.
E, assim, entende -se \textbf{que} \textbf{um} dos principais pontos da Internet of Things ( IoT ) é a ampla abertura para o uso do conceito/aplicações, quando possível.
Para reforçar esse objetivo, a maior parte dos dispositivos ( objetos ) podem ser controlados a partir de \textbf{um} único smartphone ou tablete, desde \textbf{que} esteja, é claro, conectado em \textbf{uma} rede Wi-Fi ou 3G/4G.
Alguns equipamentos permitem também conexão Bluetooth ou NFC.
A sigla NFC significa Near Field Communication, por meio dela é possível interagir com outros equipamentos, simplesmente \textbf{aproximando} o celular, sem a necessidade de fios.
A Internet das Coisas é \textbf{uma} extensão da Internet \textbf{atual}, \textbf{que} de certa forma faz com \textbf{que} objetos utilizados, diariamente, tenham capacidade computacional e de comunicação, para se conectarem à Internet.
Por meio de conexão com a rede mundial de computadores, será possível, em primeiro plano, controlar remotamente os objetos e, em segundo plano, permitir \textbf{que} os próprios objetos sejam acessados como possíveis provedores de serviços.
Os objetos comuns, isto é, do cotidiano, com novas habilidades, criam \textbf{um} diferencial de \textbf{inúmeras} oportunidades no âmbito acadêmico, no industrial, no comércio, na saúde e educação, porém, estas possibilidades apresentam riscos e acarretam amplos desafios técnicos e sociais para gerenciamento de \textbf{uma} política global de implementação da IoT.
Os objetos possuem capacidade de comunicação e processamento de dados aliados a sensores, capazes de transformar a utilidade destes objetos, ou seja, são os objetos inteligentes.
Atualmente, não só computadores estão conectados à grande rede, há também \textbf{uma} significativa quantidade de equipamentos, como TVs, Laptops, automóveis, smartphones, consoles de jogos, webcams e \textbf{uma} ampla lista de objetos \textbf{que} passam a fazer parte da mesma.
O cenário é recente, a pluralidade aumenta e previsões indicam \textbf{que} mais de 40 bilhões de dispositivos estarão conectados até 2020 ( Forbes 2014 ).
Usando a tecnologia e os recursos \textbf{funcionais} desses objetos, será possível detectar seu contexto e analisar, controlá -lo, tornar possível a troca de informações \textbf{uns} com os outros, acessar serviços da Internet e criar a interação com pessoas.
Simultaneamente, \textbf{uma} diversidade de novas interações e possibilidades de aplicações surgem cidades inteligentes ( Smart Cities ), saúde ( Healthcare ), casas inteligentes ( Smart Home ) e consideráveis desafios emergem ( regulamentações, segurança, padronizações ).
Os objetos, quando estabelecem comunicação com outros dispositivos, caracterizam o conceito de estarem em rede, assim sendo, considere \textbf{que} a IoT representa a próxima evolução da Internet, realizando \textbf{um} salto fantástico na capacidade de coletar, analisar e distribuir dados \textbf{que} podem transformar em informações, conhecimento e, nesse sentido, a escola pode realizar pesquisas a respeito do funcionamento da IoT, a respeito do funcionamento da internet das coisas em questões \textbf{que} envolvam a interação com energia e suas transformações, como as ondas de infravermelho são utilizadas para controle de equipamentos eletrônicos, a logística, envolvendo toda a distribuição da produção na agropecuária e outros, armazenamento, e como a IoT possibilita ações de sustentabilidade.
Sugere -se \textbf{que} o PPP da Escola contemple Projetos \textbf{que} incentivam os estudantes conhecer a IoT e, posteriormente desenvolvam modelos de IoT para o dia a dia em residências, indústrias de modo geral, setor agropecuário, escola e aplicação de controle de \textbf{ensino-aprendizagem}.
Os Projetos possibilitam a interação, o diálogo, a vanguarda, e, além disso, deve \textbf{focar} na ampliação de mundo mediante ensino de Ciências e sustentabilidade, por fim, refletir com os estudantes a respeito da distribuição dos recursos do mundo para aqueles \textbf{que} mais precisam deles e ajudar a entender o planeta para poder ser mais proativo e menos reativo.
A proposta do documento é \textbf{que} a escola organize grupos de estudos para ler documentos relacionados com a IoT e faça dela \textbf{um} instrumento com potencial facilitador de \textbf{ensino-aprendizagem}, motivação para ensino de Ciências, no qual a escola possa inovar as práticas de ensino, aliando com as \textbf{que} já utiliza.
As redes IoT, em cenários educacionais, proporcionam vantagens importantes para \textbf{um} desempenho \textbf{aprimorado} do ensino-aprendizagem aplicado nas escolas e universidades ao redor do mundo.
Assim, em primeiro lugar, redes IoT permitem \textbf{que} objetos ( Things ) ou pessoas sejam rastreadas e localizadas em suas respectivas instituições, reduzindo demanda e tempo de atividades, como a verificação da frequência escolar do estudante, a conexão entre atividades pedagógicas desenvolvidas por professores, o rastreio de objetos como data show, notebook, laptops, livros ou objetos do patrimônio institucional.
Em segundo lugar, abordando os aspectos cognitivos, é possível, com a utilização de dispositivos da neurociência, identificar cuidadosamente os aspectos afetivos do estudante, auxiliando o professor a propor intervenções pedagógicas e/ou conteúdo para \textbf{um} indivíduo ou para a sala como \textbf{um} todo.
Desse modo, é possível tornar eficiente e eficaz o monitoramento e motivação para o engajamento dos estudantes em \textbf{uma} dada tarefa/disciplina, permitindo \textbf{que} metodologias diferenciadas de ensino e aprendizado sejam aplicadas de acordo com as características destes estudantes.
Por último, como ações de longo prazo, o uso de sistemas embarcados conectados na Internet propicia \textbf{uma} imersão sem precedentes, tornando o ambiente mais interativo e adequado às necessidades \textbf{modernas} de aprendizado. \textbf{Aqui} se apresenta, como sugestão, \textbf{uma} visão geral da tecnologia IoT e seus benefícios para a área de educação.
Fica, como apontamento, \textbf{que} o principal objetivo é descrever a tecnologia e seu estado da arte, identificando quais componentes e serviços são adequados para prover \textbf{um} aprendizado mais efetivo.
Execução de Projetos \textbf{Uma} aula com o tema referente a características da fotossíntese poderia ser dada, apresentando, além do conhecimento \textbf{teórico}, análise de dados reais adquiridos por sensores espalhados em \textbf{uma} determinada área da escola, horto florestal, mata no entorna da escola, floresta e/ou entorno da comunidade escolar.
Assim, dados empíricos reforçariam os conceitos dados em sala de aula, além de possibilitar outras atividades pedagógicas como a interação entre os estudantes, a pesquisa, a facilitação do ensino aprendizagem.
Eficiência Operacional Nesse sentido, a IoT pode ser utilizada para conservação de energia no âmbito escolar, transporte, ambiente, segurança, conservação de alimento, distribuição de alimentos, controle de frequência escolar, etc., reduzindo custos.
Por exemplo, sensores, etiquetas RFID, câmeras e dispositivos conectados podem monitorar e vigiar o ambiente inteiro, notificando instantaneamente a gestão escolar, \textbf{uma} \textbf{central} de polícia ou ao departamento de segurança.
Estudantes especiais Com relação à inclusão, a IoT pode auxiliar deficientes físicos.
Por exemplo, sensores podem detectar usuários com algum tipo de deficiência e personalizar o objeto de conhecimento, as ações pedagógicas, como o tamanho de fontes luminosas, volume de som, chamar o professor para auxiliá -lo, localização de deficientes físicos, entre outros.
O desenvolvimento e a evolução de sistemas educacionais \textbf{que} utilizem a tecnologia IoT favorecerá o sistema tradicional de ensino, além de permitir \textbf{que} novas estratégias pedagógicas possam ser otimizadas em \textbf{um} ambiente de sala de aula.
Dentre os \textbf{inúmeros} benefícios desta aplicação, destacam -se : 1.
Os estudantes aprenderão mais rápido. 2.
O processo educacional tornar -se-á mais personalizado. 3.
Sistemas educacionais com IoT propiciarão \textbf{um} \textbf{forte} incentivo à colaboração. 4.
Redes IoT permitem \textbf{que} os dispositivos sejam controlados eficazmente. 5.
Várias \textbf{tarefas} podem ser automatizadas, como frequência de alunos, avaliações, etc.
Professores terão ferramentas adequadas para realizar suas aulas de forma eficiente. 2.2.3.
Competências Específicas de Ciências da Natureza para o Ensino Fundamental 1.
Compreender as ciências como empreendimento humano, reconhecendo \textbf{que} o conhecimento científico é provisório, cultural e histórico. 2.
Compreender conceitos fundamentais e estruturas explicativas das Ciências da Natureza, bem como dominar processos, práticas e procedimentos da investigação científica, de modo a sentir segurança no debate de questões científicas, tecnológicas e socioambientais e do mundo do trabalho. 3.
Analisar, compreender e explicar características, fenômenos e processos relativos ao mundo natural, tecnológico e social, como também as relações \textbf{que} se estabelecem entre eles, exercitando a curiosidade para fazer perguntas e buscar respostas. 4.
Avaliar aplicações e implicações políticas, socioambientais e culturais da ciência e da tecnologia e propor alternativas aos desafios do mundo \textbf{contemporâneo}, incluindo aqueles relativos ao mundo do trabalho. 5.
Construir argumentos com base em dados, evidências e informações confiáveis e negociar e defender ideias e pontos de vista \textbf{que} respeitem e promovam a consciência socioambiental e o respeito a si próprio e ao outro, acolhendo e valorizando a diversidade de indivíduos e de grupos sociais, sem preconceitos de qualquer natureza. 6.
Utilizar diferentes linguagens e tecnologias digitais de informação e comunicação para se comunicar, acessar e disseminar informações, produzir conhecimentos e resolver problemas das Ciências da Natureza de forma crítica, significativa, reflexiva e ética. 7.
Conhecer, apreciar e cuidar de si, do seu corpo e bem-estar, recorrendo aos conhecimentos das Ciências da Natureza. 8.
Agir pessoal e coletivamente com respeito, autonomia, responsabilidade, flexibilidade, resiliência e determinação, recorrendo aos conhecimentos das Ciências da Natureza para tomar decisões frente a questões científico-tecnológicas e socioambientais e a respeito da saúde individual e coletiva, com base em princípios éticos, democráticos, sustentáveis e solidários.

% matched lemmas: abordagem, ancorar, aprimorar, aproximar, aqui, atual, autor, basear, central, compromisso, conceitual, conhecido, conseguinte, contemporâneo, contudo, desdobramento, despertar, disciplina, enquanto, ensino-aprendizagem, focar, forte, funcional, fundamento, hoje, humanidade, implicar, inerente, influência, interessante, interface, intermédio, inúmero, justo, lado, mente, moderno, percepção, permear, pesquisador, que, raciocínio, tarefa, teoria, teórico, um, ver, viver
\end{textsample}
